% !PRINTED 218
\weberpage{218}

For it follows first from the second of the relations (11):
\[
\Theta^2\omega = \Theta\omega\Theta^3 = \omega\Theta^2, \quad \Theta^3\omega = \Theta\omega\Theta^2 = \omega\Theta,
\]
further from the last of (11):
\[
\Theta^3\chi = \Theta\chi\omega\Theta^3 = \chi^2\omega\Theta^2\omega\Theta^3 = \chi^2\Theta,
\]
and further:
\[
\omega\chi^2 = \chi^2\omega\chi = \chi\omega,
\]
\[
\Theta^2\chi^2 = \chi\omega\Theta^3\chi = \chi\omega\chi^2\Theta = \chi^2\omega\Theta.
\]
If now we have any three elements $\chi$, $\omega$, $\Theta$ that can be
composed according to some rule, and if in this $\chi$ is of the
third, $\omega$ of the second, $\Theta$ of the fourth degree, then it follows
from the validity of the relations (11) that the elements
$\chi^{\lambda}\omega^{\mu}\Theta^{\nu}$ form a group of degree 24, which is isomorphic
to the octahedral group. For this it only remains to show
that from these assumptions it follows that the 24 elements
$\chi^{\lambda}\omega^{\mu}\Theta^{\nu}$ are all different from one another. Let us assume
that
\[
\chi^{\lambda}\omega^{\mu}\Theta^{\nu} = \chi^{\lambda'}\omega^{\mu'}\Theta^{\nu'},
\]
then it would follow that
\[
\chi^{\lambda - \lambda'} = \omega^{\mu'}\Theta^{\nu' - \nu}\omega^{-\mu},
\]
and by (11):
\[
\chi^{\lambda - \lambda'} = \omega^{\mu' - \mu}\Theta^{\pm(\nu - \nu')}.
\]
But now it follows from (11) that $\omega\Theta$, $\omega\Theta^2$, $\omega\Theta^3$ are of the
second degree, and since $\chi$ is of the third degree, this relation
can hold only if $\lambda - \lambda'$ is divisible by 3, $\mu - \mu'$ by 2 and
$\nu - \nu'$ by 4, and hence $\chi^{\lambda} = \chi^{\lambda'}$, $\omega^{\mu} = \omega^{\mu'}$, $\Theta^{\nu} = \Theta^{\nu'}$.

By the relation following from (10),
\[
\omega = \chi\Theta\chi\Theta^2
\]
we can compose $\omega$ from $\chi$ and $\Theta$, and may therefore regard $\chi$ and
$\Theta$ as \gesperrt{generating substitutions} of the group.

The two still missing fundamental forms of the eighth and twelfth
degrees are found very easily if one first forms the Hessian
covariant of $f$.

One thus obtains the fundamental form of the eighth degree:
\weq{13}{W = x_1^8 + 14\,x_1^4x_2^4 + x_2^8,}
and if one forms from $f$ and $W$ the functional determinant,
which is again a covariant of $f$ (Bd. I, \S.\,59, 60), there
results the fundamental form of the twelfth degree:
\weq{14}{K = x_1^{12} - 33\,x_1^8x_2^4 - 33\,x_1^4x_2^8 + x_2^{12}.}
